\begin{abstract}
Power-system data are graded and released field by field, but inference risk often arises from combinations: individually uninformative fields can jointly reveal unit outputs, line flows, or load. We address this mismatch by deriving field-level risk from set-level inferability, measured by the out-of-sample accuracy of retrained attackers. The budgeted marginal risk of a field is its largest inferability gain over any background of at most $K$ other fields, relative to data already public. This quantity identifies both a field's worst-case contribution and a witness combination that explains its risk. To make the resulting audit tractable, an amortized attacker combines masked pre-training with a closed-form readout for each field set; retraining selected backgrounds then certifies lower bounds on field risk. Experiments cover ten targets across RTS-GMLC, the Australian National Electricity Market, PJM, and CAISO, with 124,425 exhaustively evaluated field sets. Single-field grading recovers only 13--17\% of critical fields and leaves over 95\% of minimal unsafe combinations intact. Scan--certify recovers up to 88\% with no false positives against the retrained reference, while estimate-driven withholding breaks 94--99\% of unsafe combinations with near-minimal withholding. Ablations identify the subset-specific readout as the main source of estimator accuracy. Increasing the background budget from two to three fields changes marginal risk little for most targets but reveals additional threshold-crossing combinations relevant to release decisions.
\end{abstract}
